\chapter{خطاهای رایج و راه حل آن‌ها}
\label{app:errors}

پیام خطا دوست شماست؛ اما در آغاز، خواندنش کمی تمرین می‌خواهد. هر پیام خطای
سلام سه بخش دارد: خود پیام، شماره‌ی خط و ستون به شکل \code{خط:ستون}، و
تکه‌ای از همان خط با یک نشانه‌ی \code{\textasciicircum} زیر جایی که مشکل
دیده شده است. در این پیوست رایج‌ترین پیام‌هایی را که تازه‌کارها می‌بینند،
همراه با علت و راه حل آورده‌ایم.

\section*{رشته‌ی پایان‌نیافته}

\begin{outputfa}
خطا: رشته‌ی پایان‌نیافته
\end{outputfa}

\textbf{علت.} نقل قول پایانی یک رشته فراموش شده، یا به جای \code{"} از
گیومه‌ی فارسی «» استفاده شده است.
\textbf{راه حل.} در خطی که پیام می‌گوید، هر رشته را نگاه کنید و مطمئن شوید
با \code{"} باز و بسته شده است.

\section*{'تمام' برای بستن بلوک انتظار می‌رفت}

\begin{outputfa}
خطا: 'تمام' برای بستن بلوک انتظار می‌رفت
\end{outputfa}

\textbf{علت.} یک بلوک (تابع، \fcode{اگر}، حلقه) با دونقطه باز شده اما
\fcode{تمام} آن نوشته نشده است. معمولاً شماره‌ی خط به پایان فایل اشاره
می‌کند، چون سلام تا آخر دنبال \fcode{تمام} گشته و پیدا نکرده است.
\textbf{راه حل.} تعداد دونقطه‌های آغاز بلوک و تعداد \fcode{تمام}‌ها را
بشمارید. تورفتگی درست، پیدا کردن جای گمشده را آسان می‌کند.

\section*{':' برای باز کردن بلوک انتظار می‌رفت}

\textbf{علت.} دونقطه‌ی پایان خط \fcode{تابع اصلی}، \fcode{اگر}، \fcode{تکرار}
یا \fcode{وگرنه} فراموش شده است.
\textbf{راه حل.} دونقطه را در پایان همان خط بگذارید.

\section*{شناسه‌ی ناشناخته}

\begin{outputfa}
خطا: شناسه‌ی ناشناخته 'سنن'
\end{outputfa}

\textbf{علت.} نامی به کار رفته که تعریف نشده است. بیشتر وقت‌ها یک غلط
املایی است (\fcode{سنن} به جای \fcode{سن})، یا متغیری که درون بلوک
دیگری تعریف شده و این‌جا در دسترس نیست، یا تابعی که بدون پرانتز صدا زده
شده است.
\textbf{راه حل.} املای نام را با تعریفش مقایسه کنید. اگر متغیر درون یک
حلقه یا تابع دیگر تعریف شده، آن را به جای مناسب منتقل کنید.

\section*{متغیر استفاده‌نشده}

\begin{outputfa}
خطا: متغیر استفاده‌نشده 'سن'
\end{outputfa}

\textbf{علت.} متغیری تعریف شده اما هیچ‌جا خوانده نشده است.
\textbf{راه حل.} یا از آن استفاده کنید، یا خط تعریفش را حذف کنید. اگر
فقط برای آزمایش موقت آن را نگه می‌دارید، می‌توانید نامش را با
\code{\_} شروع کنید تا سلام از آن بگذرد.

\section*{تابع استفاده نشده است}

\textbf{علت.} تابعی تعریف شده اما هیچ‌جا فراخوانی نشده است.
\textbf{راه حل.} آن را در جایی فراخوانی کنید یا حذفش کنید. فراموش نکنید
که فراخوانی پرانتز می‌خواهد: \fcode{خط جداکننده()}.

\section*{نمی‌توان متغیر تغییرناپذیر را دوباره مقداردهی کرد}

\textbf{علت.} به متغیری که بدون \fcode{متغیر} تعریف شده، با \code{=} مقدار
تازه داده‌اید.
\textbf{راه حل.} در خط تعریف، پیش از نام، \fcode{متغیر} بنویسید. اگر واقعاً
قرار نبوده عوض شود، خطی را که مقدار تازه می‌دهد بازبینی کنید؛ شاید
اشتباه از آن‌جاست.

\section*{کلمه‌ی رزرو شده}

\begin{outputfa}
خطا: 'بسته' یک کلمه‌ی رزرو شده است و نمی‌توان آن را به عنوان نام به کار برد
\end{outputfa}

\textbf{علت.} نام متغیر یا تابع، یا یکی از واژه‌های آن، جزو واژه‌های خود
زبان است (پیوست \ref{app:keywords}).
\textbf{راه حل.} نام دیگری انتخاب کنید یا دو واژه را به هم بچسبانید.

\section*{نمی‌توان نوع الف را به نوع ب نسبت داد}

\begin{outputfa}
خطا: نمی‌توان 'str' را به 'i32' نسبت داد
\end{outputfa}

\textbf{علت.} مقداری از یک نوع در جای مقداری از نوع دیگر گذاشته شده: رشته
در متغیر عددی، عدد اعشاری در متغیر صحیح، آرگومان با نوع اشتباه به تابع.
در این پیام‌ها \code{str} یعنی رشته، \code{i32} یعنی عدد صحیح،
\code{f64} یعنی عدد اعشاری و \code{bool} یعنی منطقی.
\textbf{راه حل.} نوع‌ها را یکسان کنید: با \fcode{بعنوان}، با
\fcode{به صحیح()} یا با اصلاح مقدار.

\section*{شرط باید بولی باشد}

\textbf{علت.} در \fcode{اگر}، \fcode{تاوقتی} یا عملگر \code{?:} چیزی نوشته
شده که مقدار منطقی نیست؛ مثلاً یک عدد.
\textbf{راه حل.} شرط را به یک مقایسه تبدیل کنید: به جای \fcode{اگر تعداد:}
بنویسید \fcode{اگر تعداد \lr{>} 0:}.

\section*{تابع ممکن است بدون بازگشت تمام شود}

\textbf{علت.} تابعی که نوع بازگشتی اعلام کرده، در یکی از مسیرها (مثلاً
وقتی هیچ شرطی برقرار نیست) به \fcode{بازگشت} نمی‌رسد.
\textbf{راه حل.} در پایان تابع یک \fcode{بازگشت} برای حالت پیش‌فرض
بگذارید.

\section*{برنامه اجرا می‌شود اما تمام نمی‌شود}

\textbf{علت.} این خطا نیست، حلقه‌ی بی‌پایان است: شرط یک \fcode{تاوقتی}
هرگز نادرست نمی‌شود.
\textbf{راه حل.} صفحه‌ی ویرایشگر را دوباره بارگذاری کنید. در حلقه به دنبال
متغیری بگردید که باید در هر دور تغییر می‌کرده و نکرده است.

\section*{برنامه اجرا می‌شود اما خروجی اشتباه است}

سخت‌ترین نوع اشکال، اشکالی است که سلام نمی‌تواند تشخیص دهد، چون برنامه
از نظر زبان درست است و فقط چیزی غیر از آنچه شما می‌خواستید انجام می‌دهد.
چند راهنمایی:
\begin{itemize}
  \item در چند جای برنامه \fcode{چاپ} موقت بگذارید و مقدار متغیرها را در
        میانه‌ی کار ببینید. بعد از پیدا کردن اشکال، آن‌ها را پاک کنید.
  \item به تقسیم عددهای صحیح مشکوک باشید (\code{7 / 2} برابر \code{3}
        است) و به ترتیب انجام عملگرها.
  \item مطمئن شوید که در شرط‌ها \code{==} نوشته‌اید، نه \code{=}.
  \item در حلقه‌ها بررسی کنید که شمارنده از کجا شروع و کجا تمام می‌شود؛
        اشتباه‌های «یکی کم یا یکی زیاد» بسیار رایج‌اند.
  \item مطمئن شوید که متغیرهای جمع‌کننده بیرون حلقه تعریف شده‌اند.
  \item برنامه را با کوچک‌ترین ورودی ممکن (یک عضو، صفر، عدد منفی)
        امتحان کنید.
\end{itemize}
